\documentclass[../../main.tex]{subfiles}
\graphicspath{{../../image/}{../../figure/}}
\begin{document}

\section{深度网络}\label{section:ch06-深度网络}

我们通过使用可学习参数控制线性回归/分类模型的基函数，阐述了发展神经网络的动机，从而产生了图~\ref{fig:ch06-mlp}~所示的两层网络。多年来它一直是最广为使用的架构——主要因为很难有效训练超过两层的网络。然而最近训练技术的发展使多层网络变得有效（参见第~7~章）。

注意此类网络中"层数"的术语存在一些混乱：图~\ref{fig:ch06-mlp}~所示网络有时被描述为"三层网络"（对单元层数计数，把输入视为单元），有时被描述为"单隐藏层网络"（对隐藏单元层数计数）。我们建议将图~\ref{fig:ch06-mlp}~所示网络称为"\textbf{两层网络}"，因为可学习权重的层数对网络性质十分重要。

我们可以轻松地把两层网络架构~\eqref{align:ch06-mlp-scalar}~扩展到任意有限层数 $L$，其中第 $l=1,\cdots,L$ 层计算
\begin{align}
\mathbf{z}^{(l)} = h^{(l)}\!\left(\mathbf{W}^{(l)}\mathbf{z}^{(l-1)}+\mathbf{b}^{(l)}\right),
\label{align:ch06-deep-layer}
\end{align}
其中 $h^{(l)}$ 是与第 $l$ 层相关联的激活函数，$\mathbf{W}^{(l)}$ 表示相应的权重和偏置参数矩阵。此外 $\mathbf{z}^{(0)} = \mathbf{x}$ 表示输入向量，$\mathbf{z}^{(L)} = \mathbf{y}$ 表示输出向量。

虽然图~\ref{fig:ch06-mlp}~所示的两层网络具有通用近似能力，但具有超过两层的网络有时可以用比两层网络少得多的参数来表示给定函数。Montufar 等人（2014）指出，网络函数将输入空间划分成多个区域，这些区域在网络深度上是指数级增长的，但在隐藏层宽度上却是多项式级的；使用两层网络表示相同函数则需要指数级数量的隐藏单元。

\subsection{层次化表示}\label{section:ch06-层次化表示}

虽然通用近似定理是有趣的结果，但另一个更令人信服的探索深度网络的理由是：网络架构编码了一种特定形式的归纳偏置，即输出空间与输入空间之间通过层次化表示相关联。这方面一个很好的例子是图像中物体识别（参见第~10~章）：像素与高级概念（如"猫"）之间的关系非常复杂且非线性，对两层网络而言是极具挑战性的识别任务；然而深度网络会首先在浅层检测低级特征（如边缘），然后在后续层中将它们结合起来以得到更高级的特征（如眼睛或胡须），之后继续在深层组合这些特征以检测猫的存在。

这可以看作一种\textbf{组合}（compositional）归纳偏置：更高级对象（如猫）由低级对象（如眼睛）组成，而这些低级对象又由更低级元素（如边缘）组成。也可以从相反角度考虑图像生成：从低级特征（如边缘）开始，组合起来形成简单形状（如圆），之后进一步组合形成更高级对象（如猫）。每个阶段都有许多不同方式组合不同组件，随着深度增加，可能性数量呈指数级增长。

\subsection{分布式表示}\label{section:ch06-分布式表示}

在概念上，隐藏层中的每个单元都可以视为表示网络在该层的一个"特征"：激活值高表示相应特征存在，激活值低表示不存在。当给定层中有 $M$ 个单元时，网络可以表示 $M$ 个不同特征。然而网络也可能学习到不同表示——其中隐藏单元的\emph{组合}表示特征，从而使具有 $M$ 个单元的隐藏层能表示 $2^M$ 个不同特征。

例如考虑一个用于处理人脸图像的网络：每张特定人脸图像可能有（或没有）眼镜、有（或没有）帽子、有（或没有）胡须，从而导致 $2^3=8$ 种组合。尽管这可以用 $8$ 个单元来表示（检测到相应组合时对应单元"打开"），但也可以更紧凑地只使用 $3$ 个单元，每个单元代表一个属性。这些属性可以独立存在（尽管统计上它们的存在可能一定程度上相关）。稍后我们将在第~10~章详细探讨深度学习网络在训练过程中自行发现的内部表示类型。

\subsection{表示学习}\label{section:ch06-表示学习}

我们可以把连续经过深度网络不同层视为对数据进行变换，以便更容易地完成任务。例如成功学习将皮肤病变分类为良性或恶性（参见~1.1.1~小节）的网络必须学会将原始图像数据转换为由最终隐藏层输出表示的新空间，以便网络最终层可以区分这两类数据。这个最终层可以视为一个简单线性分类器，因此在最后隐藏层的表示中，两类数据必须能被一个线性面很好地分开。

这种发现数据的非线性变换使后续任务更容易完成的能力称为\textbf{表示学习}（representation learning）~\citep{BengioCourvilleVincent2012}。学习到的表示有时称为\textbf{嵌入空间}（embedding space），它们由网络中一个隐藏层的输出给出，从而对于任何输入向量（无论来自训练集还是新数据集）都可以通过网络前向传播转换为这个表示。

表示学习异常强大的原因在于它使我们能够利用无标签数据。通常收集大量无标签数据很容易，但获取相关标签可能更困难：例如车辆行驶在城市周边时，车载摄像头可以收集大量城市场景图像，但识别相关物体（如行人、路标）却需要成本高昂且耗时的人工标注。

从无标签数据中学习称为\textbf{无监督学习}（unsupervised learning），人们已开发出许多不同算法。例如可以训练神经网络将图像作为输入并创建相同图像作为输出；为使其成为一项具有挑战性的任务，网络可能使用比图像像素数更少的隐藏层单元，从而要求网络学习对图像进行某种压缩。这只需要无标签数据，因为训练集中每张图像同时充当输入向量和目标向量。这种网络又称为\textbf{自编码器}（参见~19.1~节），训练目标是使网络发现对解决其他任务（如图像分类）有用的数据内部表示。

无监督学习在成功训练第一个深度网络（卷积网络除外）方面扮演了重要角色：网络每一层首先使用无监督学习进行预训练，然后整个网络再使用基于梯度的监督训练方式。后来人们发现预训练阶段可以省略，在适当条件下可以纯粹使用监督学习从头训练深度网络。然而预训练和表示学习仍然是深度学习的核心。最值得一提的预训练是在自然语言处理中进行的（参见第~12~章）：Transformer 模型在大量文本上进行训练，从而学习到高度复杂的语言内部表示，使其各项能力达到甚至超越人类水平。

\subsection{迁移学习}\label{section:ch06-迁移学习}

从一个特定任务学习到的内部表示也可能对相关任务有帮助。例如在一个大型有标签常见对象数据集上训练的网络，可以学习到如何将图像表示转换为更适合对象分类的表示；然后可以使用一个较小的有标签数据集对网络最终分类层重新训练，以创建一个病变分类器（参见~1.1.1~小节）。这是\textbf{迁移学习}（transfer learning）~\citep{Hospedales2021}的一个示例，相较于仅使用病变图像数据训练可以获得更高准确率，因为网络可以利用自然图像的通用共性。

一般来说，迁移学习可以通过使用来自相关任务 $B$ 的数据来提升任务 $A$ 的性能，其中任务 $A$ 训练数据匮乏而任务 $B$ 训练数据更丰富。两个任务应具有相同类型输入，且任务之间应存在一些共性，使得从任务 $B$ 学到的低级特征或内部表示对任务 $A$ 有帮助。当研究卷积网络（参见第~10~章）时，我们会发现许多图像处理任务在网络早期层中需要相似低级特征，而后续层则专门用于特定任务，因此这样的网络非常适合迁移学习。

当任务 $A$ 训练数据非常稀缺时，可能只需重新训练网络最终层；如果训练数据更多，重新训练多个层也可行。使用一个任务学习参数，然后应用于一个或多个其他任务的过程称为\textbf{预训练}（pre-training）。对于新任务，相较于对整个网络应用随机梯度下降，更高效的方式是将新训练数据一次通过参数固定的预训练网络，以便在新表示中评估训练输入，然后在仅包含最终层的较小网络上应用基于迭代梯度的优化。除了把预训练网络作为固定预处理器之外，也可以应用\textbf{微调}（fine-tuning），即将整个网络适配到任务 $A$ 的训练数据上；这通常通过非常小的学习率和有限次迭代完成，以确保网络不会过拟合新任务中相对较小的可用数据集。

跨多个任务学习可以扩展到\textbf{多任务学习}（multitask learning）~\citep{Caruana1997}：网络同时联合学习多个相关任务。例如希望建立一个垃圾邮件过滤器，允许不同用户有不同分类器以适应各自偏好；训练数据可能包括许多不同用户的垃圾邮件和非垃圾邮件样本，但任何一个用户的样本数量可能相当有限，因此为每个用户单独训练分类器可能效果不佳。相反，可以将数据集组合并训练一个单一较大网络，该网络共享早期层，但对不同用户后续层具有单独可学习参数。跨任务共享数据允许网络利用任务之间的共性，从而提高对所有用户的准确率。大量训练样本使我们能够使用具有更多参数的更深层网络，进一步提高性能。

跨多个任务学习还可以扩展到\textbf{元学习}（meta-learning），也称"学会学习"（learning to learn）~\citep{Hospedales2021}。多任务学习旨在对一组固定任务进行预测，而元学习的目标是对训练期间未见过的未来任务进行预测。这不仅可以通过跨任务学习共享内部表示来实现，还可以通过学习学习算法本身来实现。例如当新类的有标签样本非常少时，元学习可以用于促进分类模型泛化到新类上，称为\textbf{小样本学习}（few-shot learning）；当只使用一个有标签样本时，则称为\textbf{单样本学习}（one-shot learning）。

\subsection{对比学习}\label{section:ch06-对比学习}

\textbf{对比学习}（contrastive learning）是最常见和强大的表示学习方法之一~\citep{GutmannHyvarinen2010,OordLiVinyals2018,ChenKornblith2020}。它背后的思想是学习一种表示，使得在嵌入空间中某些输入对（称为\textbf{正对}）之间距离近，而其他输入对（称为\textbf{负对}）之间距离远。直觉上，如果选择语义相似的为正对、语义不相似的为负对，那么将学习到一个表示空间，在该空间中相似输入之间距离近，从而使完成下游任务（如分类）更容易。

与其他形式的表示学习一样，我们通常不直接使用训练好的网络输出，而是使用某个早期层的激活来构建嵌入空间。对比学习任务与其他大多数机器学习任务不同：给定输入的误差函数仅取决于其他输入，而非每个输入单独的标签或目标输出。

给定数据点 $\mathbf{x}$（称为\textbf{锚点} anchor），指定另一个数据点 $\mathbf{x}^{+}$ 与之组成一个正对；还必须指定一组数据点 $\{\mathbf{x}^{-}_{1},\cdots,\mathbf{x}^{-}_{N}\}$，其中每个点都与 $\mathbf{x}$ 组成一个负对。我们需要一个损失函数：如果 $\mathbf{x}$ 和 $\mathbf{x}^{+}$ 的表示较为接近则给予奖励，同时拉大 $\mathbf{x}$ 与每个 $\mathbf{x}^{-}$ 之间的距离。

一个示例，也是对比学习中最常用的损失函数，称为 InfoNCE 损失函数~\citep{GutmannHyvarinen2010,OordLiVinyals2018}，其中 NCE 表示"噪声对比估计"（noise contrastive estimation）。假设有一个神经网络函数 $f_{\mathbf{w}}(\mathbf{x})$，用于将来自输入空间的点 $\mathbf{x}$ 映射到一个由可学习参数 $\mathbf{w}$ 控制的表示空间。该表示是归一化的，即 $\|f_{\mathbf{w}}(\mathbf{x})\|=1$。对于数据点 $\mathbf{x}$，InfoNCE 损失函数定义如下：
\begin{align}
\mathcal{L}_{\mathrm{NCE}} = -\ln\!\left(\frac{\exp(f_{\mathbf{w}}(\mathbf{x})\cdot f_{\mathbf{w}}(\mathbf{x}^{+})/\tau)}{\exp(f_{\mathbf{w}}(\mathbf{x})\cdot f_{\mathbf{w}}(\mathbf{x}^{+})/\tau)+\sum_{j=1}^{N}\exp(f_{\mathbf{w}}(\mathbf{x})\cdot f_{\mathbf{w}}(\mathbf{x}^{-}_{j})/\tau)}\right).
\label{align:ch06-infonce}
\end{align}

在这个函数中，锚点表示 $f_{\mathbf{w}}(\mathbf{x})$ 和正样本表示 $f_{\mathbf{w}}(\mathbf{x}^{+})$ 之间的余弦相似度给出了在习得空间中度量正对样本有多近的方法；同样的方法也可用于评估锚点与负对样本之间的接近程度。注意这个函数类似于分类交叉熵误差函数，其中正对的余弦相似度给出了标签类的对数概率，而负对的余弦相似度给出了错误类的对数概率。负对非常重要：如果没有它们，嵌入空间将简单地学习将每个点映射到相同表示的退化解。

一个特定的对比学习算法主要由正对和负对的选择方式定义，也就是如何使用先验知识来指定什么是好的表示。例如考虑学习图像表示：一个常见选择是通过\emph{破坏}输入图像来创建正对，这种方式保留了图像的语义信息，同时极大地改变了像素空间中的图像~\citep{Wu2018,He2019,ChenKornblith2020}；破坏可以是旋转、平移和颜色转换。数据集中其他图像可用来创建负对。这种对比学习方法称为\textbf{个体判别}（instance discrimination）。

然而如果可以访问类标签，则可以使用同一类别的图像作为正对，而使用不同类别的图像作为负对。这放宽了对指定表征应保持不变的增强方法的依赖，并且避免了将两个语义相似图像视为负对。这称为\textbf{监督对比学习}（supervised contrastive learning）~\citep{Khosla2020}，因为它依赖于类标签，并且通常能得到比仅使用交叉熵分类学习的表示更好的结果。

正对和负对的成员不一定非要来自同一数据模态。在对比语言--图像预训练（Contrastive-Language Image Pretraining, CLIP）~\citep{Radford2021}中，一个正对由一幅图像和对应的文本说明组成；用两个单独函数将每个模态输入分别映射到相同表示空间，其中 $f$ 表示从图像到表示空间的映射，$g$ 表示从文本输入到表示空间的映射。还需要从数据集中获取一组图像 $\{\mathbf{x}_1,\cdots,\mathbf{x}_N\}$，并假设对它们来说文本说明 $\mathbf{y}^{+}$ 是不恰当的，还需一组同样不匹配输入图像 $\mathbf{x}$ 的文本说明 $\{\mathbf{y}^{-}_{1},\cdots,\mathbf{y}^{-}_{N}\}$。损失函数由下式给出：
\begin{align}
\mathcal{L}_{\mathrm{CLIP}} = -\ln\!\left(\frac{\exp(f(\mathbf{x}^{+})\cdot g(\mathbf{y}^{+})/\tau)}{\sum_{j=0}^{N}\exp(f(\mathbf{x}^{+})\cdot g(\mathbf{y}^{-}_{j})/\tau)}\right)
-\ln\!\left(\frac{\exp(g(\mathbf{y}^{+})\cdot f(\mathbf{x}^{+})/\tau)}{\sum_{j=0}^{N}\exp(g(\mathbf{y}^{+})\cdot f(\mathbf{x}^{-}_{j})/\tau)}\right),
\label{align:ch06-clip}
\end{align}
其中 $\mathbf{x}^{+}$ 和 $\mathbf{y}^{+}$ 表示一个正对。损失函数中的两项确保了：(a) 图像表示与对应的文本说明表示相对于其他图像表示而言是接近的；(b) 文本说明表示与它所描述的图像表示相对于其他文本说明表示而言是接近的。虽然 CLIP 使用文本--图像对，但任何具有配对模态的数据集都可以用来学习表示。这通常称为\textbf{弱监督学习}，因为它依赖于带文本说明的图像，这些数据通常从互联网上爬取，比手动标注图像类别更容易。

\subsection{通用网络结构}\label{section:ch06-通用网络结构}

到目前为止我们探讨了组织成一系列全连接层的神经网络架构。然而由于网络图与其数学函数之间存在直接对应关系，我们可以通过考虑更复杂的网络图来设计更通用的网络映射。这些网络必须限制为\textbf{前馈}（feed forward）网络，换句话说不能有闭合的有向环，以确保输出是输入的确定性函数。在这样的网络中，每个（隐藏或输出）单元计算的函数为
\begin{align}
a_k = \sum_{j\in \mathrm{pa}(k)} w_{kj} z_j + b_k, \qquad z_k = h\!\left(a_k\right),
\label{align:ch06-graph-net}
\end{align}
其中 $\mathrm{pa}(k)$ 表示节点 $k$ 的祖先（ancestor）集，换句话说就是指向单元 $k$ 的单元所构成的集合，而 $\{b_k\}$ 表示相关联的偏置参数。对于应用于网络输入的一组给定值，连续应用~\eqref{align:ch06-graph-net}~可以计算网络中所有单元的激活，包括输出单元的激活。

\subsection{张量}\label{section:ch06-张量}

可以看到线性代数在神经网络中发挥着核心作用，其中数据集、激活和网络参数等表示为标量、向量和矩阵。此外我们还遇到了更高维度的变量。例如考虑一个包含 $N$ 张彩色图像的数据集，每张图像高度为 $I$ 像素，宽度为 $J$ 像素。每个像素可通过其在图像内所处的行和列索引，并具有红、绿、蓝值；对数据集中每张图像我们都有这样一个值，因此可以用一个四维数组 $\mathcal{X}$ 来表示：
\begin{align}
x_{nij k}, \qquad i\in\{1,\cdots,I\},\ j\in\{1,\cdots,J\},\ k\in\{1,2,3\},\ n\in\{1,\cdots,N\},
\end{align}
其中 $i,j$ 是图像内行和列的索引，$k\in\{1,2,3\}$ 是红、绿、蓝强度索引，$n\in\{1,\cdots,N\}$ 是数据集中特定图像的索引。

这些高维数组称为\textbf{张量}（tensor），而标量、向量和矩阵是张量的特例。等到后面讨论更复杂的神经网络架构时，我们会看到许多张量的例子。像 GPU 这样的大规模并行处理器特别适合处理张量。

\end{document}
